\subsection{Index residency and construction workspace}

Compression determines whether an index can remain on the GPU, but
construction has its own capacity constraint. BGE occupies 11.93\,GiB
during search yet reaches a sampled 26.08\,GiB during construction.
Assignment buffers cause this peak before the gather buffers are
allocated. The two-pass build limits gather storage; it does not remove
all other temporary allocations. Feasibility therefore requires both
a resident-byte model and a phase-by-phase allocation model
(\cref{sec:4-5,sec:7-7}).

\subsection{Codebook structure and lookup cost}

The smallest table is not the fastest. Four- and eight-way splitting
store 16 entries per subspace, compared with 32 for the two-way split,
but need more lookups and additions. Two tables are fastest in 22 of
24 settings (\cref{sec:7-3}). Likewise, removing the table altogether
saves shared memory but loses speed in the earlier tests summarised in
\cref{tab:8-stopped}. The useful comparison is therefore reconstruction quality at the actual
GPU scoring cost.

\subsection{Parallelism must repay selection and merging}

Small batches leave too few independent queries to fill the GPU.
Splitting each query across blocks supplies more work, but repeats table
staging and local selection and adds a merge. Eight groups gain
$2.06$--$2.76\times$ at batch 8; at batch 1,000, independent-bound
grouping falls to $0.46\times$ one-block throughput in the profiled
high-budget case (\cref{sec:7-6}). This does not mean that large batches
should always use one block: grouping also handles uneven workloads,
and supplies Vogue's lead at recall 0.97. Its 0.895 throughput ratio on
a low-probe setting shows the rule's limit. The spread of candidate
counts does not measure their total volume; a better scheduler would
price scan, selection and merge work together.

\subsection{Selection correctness and search quality}

Exact selection of computed scores is a separate requirement from
finding the nearest original vectors. The grouped selector's proof and
CPU sorting checks establish that parallel decomposition retains the
specified primary shortlist (\cref{sec:5-4,sec:7-strict-selection}).
They do not remove IVF routing losses or quantisation error. Testing
these layers separately distinguishes an implementation defect from a
limitation of the representation or refinement budget.

\subsection{Refinement speed and budget savings}

How much a smaller budget saves depends on the rest of the search.
Let $T_0>0$ be fixed non-refinement time, $F_h\geq F_\ell>0$ the
refinement times at high and low budgets, and $\eta\in(0,1]$ a common
multiplier that makes both refinement kernels faster. The low-budget
speedup and its derivatives are
\begin{align}
 S(\eta,T_0)&=\frac{T_0+\eta F_h}{T_0+\eta F_\ell},
 \label{eq:8-refinement-sensitivity}\\
 \frac{\partial S}{\partial\eta}
 &=\frac{T_0(F_h-F_\ell)}{(T_0+\eta F_\ell)^2}\geq0,\qquad
 \frac{\partial S}{\partial T_0}
 =\frac{\eta(F_\ell-F_h)}{(T_0+\eta F_\ell)^2}\leq0.\nonumber
\end{align}
For $F_h>F_\ell$, reducing $\eta$ reduces the relative benefit of
cutting the budget, whereas reducing $T_0$ increases it. This model
explains why kernel optimisation and budget reduction cannot be treated
as independent multiplicative gains. It assumes the same refinement
speedup at both budgets and fixed non-refinement work. Actual scan
maintenance can change with the budget, so adoption still depends on
measured total time and quality (\cref{sec:7-5,sec:7-stream-cost}).
